\chapter{Experimentos detallados}\label{ap:experimentos}

Este apéndice reúne los experimentos que sustentan las decisiones del capítulo~\ref{cap:metodologia}. Salvo que se diga
otra cosa, las cifras son medidas F contra las referencias corregidas de oído, con las pistas de Demucs de semilla 0.

\section{Altura de la línea: cambios probados y descartados}

Cada cambio se midió con una réplica de la celda de la línea fuera del cuaderno, que reproduce la F guardada con
$\pm0.9$ puntos en la voz aguda y exactamente en el bajo. Solo se adoptaba si subía en las tres canciones
(tabla~\ref{tab:ronda-altura}).

\begin{table}[ht]
  \centering\small
  \caption{Cambio en puntos de F (Sarniezz / Mata Zyklek / Fabienk).}\label{tab:ronda-altura}
  \begin{tabularx}{\textwidth}{@{}>{\raggedright}X l@{}}
    \toprule
    Cambio & Resultado \\
    \midrule
    Voto entre vueltas para el cuarto de tono & arregla 1 / 17 / 5 errores, rompe 6 / 4 / 0 aciertos \\
    Altura fina por el pico de cada armónico & $-0.1$ / $-0.2$ / $+0.1$ \\
    Bajar una octava cuando suenan los armónicos de la nota grave & $-0.4$ / $+0.2$ / $+0.6$ \\
    Piso de la voz aguda en MIDI 45 (en lugar de 47) & $-2.9$ / $-1.6$ / $-1.6$ \\
    Atracción al registro más ancha o local & $-0.3$ a $-0.9$ / $-2.7$ a $-3.8$ / $+0.2$ a $+0.5$ \\
    Umbral de ataques de la guitarra 0.2 (en lugar de 0.3) & $-0.4$ / $0.0$ / $-0.2$ \\
    Nota del bajo cuando la guitarra lo dobla & $-0.5$ / $-2.4$ / $-0.2$ \\
    Bajo: ataques de banda completa & $+0.5$ / $+2.3$ / $-1.1$ \\
    \bottomrule
  \end{tabularx}
\end{table}

Los errores de cuarto de tono se repiten igual en todas las vueltas. En el 74--83\,\% de ellos, el pico de la pista de
guitarra está más cerca de la altura que leyó la herramienta que de la partitura (mediana: 34--45 cents desde la
partitura), y se concentran en pocas notas: apuntan a la referencia, lo que motivó el arbitraje a ciegas. En Fabienk, el
10\,\% de la melodía está por debajo de MIDI 47, donde la guitarra dobla al bajo; ahí la voz aguda toma un armónico, y bajar
el piso hace caer la línea una octava en todas las canciones. En el bajo, el error mayor son los re-ataques perdidos
(9--11\,\% en Mata Zyklek y Fabienk).

\section{El peso de la capa más nueva}

Se probó dar más peso a la capa más nueva del looper en la voz aguda, con tres valores del parámetro de edad. La melodía
cambió, pero la fundamental del acorde no, salvo en un valor extremo para Mata Zyklek (tabla~\ref{tab:edad}). Estas
corridas son de una versión anterior del cuaderno.

\begin{table}[ht]
  \centering\small
  \caption{Melodía principal y fundamental con tres pesos de edad (0.5, 0.7 y 0.85).}\label{tab:edad}
  \begin{tabular}{@{}lll@{}}
    \toprule
    Canción & Melodía principal & Fundamental \\
    \midrule
    Sarniezz & 46.0 · 46.8 · 47.9 & 72.3 · 72.3 · 72.3 \\
    Fabienk & 40.8 · 42.8 · 44.2 & 82.3 · 82.3 · 82.3 \\
    Mata Zyklek & 60.4 · 62.8 · 65.3 & 90.9 · 98.1 · 98.1 \\
    \bottomrule
  \end{tabular}
\end{table}

\section{Separación por sección}

La tabla~\ref{tab:sep-seccion} detalla, sección por sección de Sarniezz, la melodía principal con cada separador
($n$: número de notas de la referencia).

\begin{table}[ht]
  \centering\small
  \caption{Melodía principal de Sarniezz por sección: Demucs frente a BS-RoFormer SW.}\label{tab:sep-seccion}
  \begin{tabular}{@{}lrrr@{}}
    \toprule
    Sección & $n$ & Demucs & RoFormer \\
    \midrule
    Riff del bajo (c.~9--13) & 2 & 28.6 & 28.6 \\
    Riff, puente (c.~14) & 9 & 35.3 & 33.3 \\
    Riff en bajo y guitarra 2 (c.~15--18) & 96 & 82.2 & 84.1 \\
    Riff y acorde de la guitarra 1 (c.~19--22) & 96 & 81.8 & 83.1 \\
    Frase 1 de la guitarra 1 sobre el riff (c.~23--26) & 86 & 12.4 & 7.5 \\
    Frase 1 con bajo suelto (c.~27--30) & 160 & 55.9 & 40.3 \\
    Frase 1 en 4/4 (c.~31--36) & 80 & 62.4 & 60.5 \\
    Frase 1 con la guitarra 2 una octava abajo (c.~37--42) & 160 & 7.6 & 10.9 \\
    Frase 1 y frase 2 aguda de la guitarra 2 (c.~43--51) & 164 & 19.3 & 43.0 \\
    Patrón del bajo en la guitarra 1 (c.~64--69) & 192 & 59.5 & 54.6 \\
    Final libre de la guitarra 1 (c.~77--82) & 24 & 45.3 & 40.0 \\
    \bottomrule
  \end{tabular}
\end{table}

\section{Columnas melódicas y patrones frente a la partitura}

La misma regla de columnas aplicada a la partitura permite medir qué tan estable es frente a los errores de transcripción
(tabla~\ref{tab:columnas}). La nota de la partitura suele estar entre las columnas de pulso de la herramienta aunque no sea
la principal. La partitura confirma además que la guitarra dobla al bajo en unísono u octava buena parte del tiempo
(38--73\,\% de las columnas), y que el resto se apoya en quintas y terceras.

\begin{table}[ht]
  \centering\small
  \caption{Columnas melódicas: coincidencia con la partitura.}\label{tab:columnas}
  \begin{tabular}{@{}lrrrrr@{}}
    \toprule
    Canción & \shortstack{Misma nota\\principal} & \shortstack{Entre las\\columnas} & \shortstack{Mismo\\bajo} &
    \shortstack{Misma nota\\del acorde} & \shortstack{Mismo\\arco} \\
    \midrule
    Sarniezz & 31\,\% & 57\,\% & 33\,\% & 32\,\% & 59\,\% \\
    Mata Zyklek & 63\,\% & 83\,\% & 81\,\% & 64\,\% & 53\,\% \\
    Fabienk & 37\,\% & 68\,\% & 29\,\% & 36\,\% & 35\,\% \\
    \bottomrule
  \end{tabular}
\end{table}

Para los patrones de ciclo, la tabla~\ref{tab:patrones} compara el patrón de cada ciclo de la herramienta con el de la
partitura, y la fracción de notas de cada vuelta que se acoge al patrón.

\begin{table}[ht]
  \centering\small
  \caption{Patrones de ciclo frente a la partitura.}\label{tab:patrones}
  \begin{tabular}{@{}lrrr@{}}
    \toprule
    Canción & Parecido de los patrones (F) & Se acogen: herramienta & Se acogen: partitura \\
    \midrule
    Sarniezz & 52\,\% & 57\,\% & 62\,\% \\
    Mata Zyklek & 86\,\% & 56\,\% & 60\,\% \\
    Fabienk & 78\,\% & 72\,\% & 67\,\% \\
    \bottomrule
  \end{tabular}
\end{table}

\section{Ruido de muestreo en los cambios de acorde}

Con pocos cambios de fundamental, las distancias entre repartos de transiciones son sobre todo ruido. Para medirlo se
sacaron dos muestras de $n$ cambios de una misma canción y se midió cuánto difieren (tabla~\ref{tab:ruido}). Con 6
cambios, como en Utzp y Yor Zarad, el ruido supera la distancia típica entre un audio y su partitura (alrededor de 0.1).
Por eso, en la distancia combinada entre círculos, las transiciones y los movimientos pesan según el número de cambios:
nada con menos de 5, $n/20$ entre 5 y 20, y completo desde 20.

\begin{table}[ht]
  \centering\small
  \caption{Distancia entre dos muestras de una misma canción según el número de cambios.}\label{tab:ruido}
  \begin{tabular}{@{}rrr@{}}
    \toprule
    Cambios & Transiciones & Movimientos \\
    \midrule
    6 & 0.28 & 0.12 \\
    10 & 0.19 & 0.07 \\
    20 & 0.11 & 0.03 \\
    40 & 0.06 & 0.02 \\
    \bottomrule
  \end{tabular}
\end{table}

\section{Las referencias de Guitar Pro}

La referencia completa de acordes se armó con las tablaturas de los PDF por pista (alturas, cuerdas, trastes, ligados y
bends) y con los audios que Guitar Pro exporta por pista (tiempos). Entre el 82 y el 88\,\% de las notas tiene su ataque
medido en el audio; el resto se interpola entre sus vecinas. El render se alinea con la grabación por DTW sobre un croma
de cuartos de tono. De las notas de la referencia corregida de oído, también están en Guitar Pro (misma altura, $\pm75$~ms)
el 92\,\% en Mata Zyklek, el 70\,\% en Fabienk y el 57\,\% en Sarniezz; sin contar la octava, el 96, el 80 y el 77\,\%. El
acuerdo entre referencias en la fundamental por compás es del 100, 92.6 y 97.0\,\% entre la corregida y la completa, y del
95.4, 49.5 y 75.5\,\% entre la corregida y Guitar Pro sola.

\section{Elección de los fragmentos de las figuras}

\paragraph{El círculo en movimiento.} El fragmento se eligió cuadro a cuadro contra la referencia corregida
(tabla~\ref{tab:fragmento}). El primero tenía más cuartos de tono, pero la mitad eran errores ya descartados por el
arbitraje a ciegas (el Fa leído como Mi+ y el Re como Re+). En el elegido, los cuartos de tono son reales y, cuando la
herramienta falla, se queda corta.

\begin{table}[ht]
  \centering\small
  \caption{Fragmentos de Mata Zyklek para el círculo en movimiento.}\label{tab:fragmento}
  \begin{tabular}{@{}lrrl@{}}
    \toprule
    Fragmento & Voz aguda & Voz grave & Cuartos de tono (referencia · círculo) \\
    \midrule
    162--192 s & 69.0\,\% & 86.8\,\% & 25\,\% · 38\,\% \\
    42--72 s (elegido) & 86.9\,\% & 94.4\,\% & 23\,\% · 17\,\% \\
    \bottomrule
  \end{tabular}
\end{table}

\paragraph{El tabroll de Fabienk.} La página de la figura~\ref{fig:tabroll} se eligió comparando las páginas posteriores a
3:00 con la partitura corregida (tabla~\ref{tab:paginas}). La página 25 marcaba cuartos de tono que la partitura no tiene;
la 32 tiene la melodía principal y casi todos sus cuartos de tono están en la partitura.

\begin{table}[ht]
  \centering\small
  \caption{Páginas candidatas del tabroll de Fabienk.}\label{tab:paginas}
  \begin{tabular}{@{}llrrr@{}}
    \toprule
    Página & Tramo & Misma nota que la melodía & Notas en alguna guitarra & \cuarto{} de tono en la partitura \\
    \midrule
    25 & 3:43--3:53 & 73\,\% & 34 de 48 & 4 de 11 \\
    32 & 4:42--4:52 & 78\,\% & 53 de 59 & 6 de 8 \\
    33 & 4:52--5:01 & 80\,\% & 48 de 52 & 3 de 4 \\
    \bottomrule
  \end{tabular}
\end{table}
